\section{The partition function and the highest changing genus}
\label{sec:partition}

Write $h=g(Y)$, $b=h-1=2(g-1)$, $y=e^{iv}$ and
$z_j=2-y^j-y^{-j}$, so that $z=z_1=(2\sin(v/2))^2$.
For a box $a=(i,j)$ in a partition $\lambda$, let $\ell(a)$ be its
hook length and put $c(\lambda)=\sum_{(i,j)\in\lambda}(j-i)$.
For $p\ge1$ define
\begin{equation}\label{eq:positiveseries}
 Z_p(Q,y)=\sum_{\lambda}(-1)^{p|\lambda|} Q^{|\lambda|}
 y^{p c(\lambda)}\prod_{a\in\lambda}z_{\ell(a)}^{\,p}.
\end{equation}
The empty partition supplies the constant term.

\begin{proposition}\label{prop:evaluation}
For the character path of Theorem~\ref{thm:character} and each $e\ge1$ the connected contribution on $Y$ is
\[
 L^Y_e(B_1;v)=[Q^e]\log Z_b(Q,y).
\]
Put $c=g-1$. On $C$ the
positive and negative endpoints have disconnected series
\[
 Z_c(q,y),\qquad Z_c(q,y)\,Z_b(q^2,y)^{-1}
 \pmod {q^{M+1}}
\]
at the endpoints through degree $M$ specified in that theorem.
\end{proposition}

\begin{proof}
The flat frame upstairs identifies the two lines with
$L_1=\mathcal O_Y$ and $L_2=K_Y$. At $t=1$, the order-zero term in
$B_1$ sends each summand to the other by a conjugate-linear map.
Scale this term by $\tau\in[0,1]$. For every $\tau$ the resulting
real operator commutes with the circle of weights $(1,-1)$, hence
with the complex structure $J'=(i,-i)$ on sections and on
bundle-valued forms.

Let $\pi$ and $a$ be the universal curve and map over
$\mathcal M=\overline{\mathcal M}_k(Y,e)$ and set
$W_i=-R\pi_*a^*L_i$ in complex $K$-theory. Put $s=k-1-be$.
Riemann--Roch gives $r_1=\operatorname{rk}W_1=k-1$,
$r_2=\operatorname{rk}W_2=k-1-2be$ and
$r_1+r_2=2s=\dim^{\mathrm{vir}}_{\CC}\mathcal M$.
At $\tau=1$ the pullback operators are injective. Their cokernels
form a bundle $V$ with
\begin{equation}\label{eq:massclass}
 [V,J']=W_1+\overline{W_2}.
\end{equation}
Indeed, the entire homotopy consists of complex-linear Fredholm maps
for $J'$, including the matching conditions at the nodes. In a nodal
chart one stabilizes the family by forms on the thick part, taking
their complex span for $J'$ and their translates under the finite
stabilizer. Enlarging the stabilization preserves the index. This
gives a single family index over $\mathcal M\times[0,1]$ by the nodal
family construction of \cite[\S2.4]{zinger}. At $\tau=0$ its negative
is the sum of the first Dolbeault cokernel index and the complex
conjugate of the second. The Dolbeault resolution in families of nodal
curves identifies these indices with the classes in
\eqref{eq:massclass} \cite[\S4]{siebertnodal}.

At $\tau=0$ the change from the ordinary complex determinant
orientation to the $J'$ orientation conjugates both cohomology groups
of the second line. The relative sign is
$(-1)^{h^0(a^*L_2)+h^1(a^*L_2)}=(-1)^{r_2}$. Transporting both
orientations along the homotopy gives
\[
 e_{\mathrm{det}}(V)=(-1)^{r_2}c_{2s}(W_1+\overline{W_2})
      =[u^0]\,e_u(W_1)e_{-u}(W_2).
\]
Here $e_u(W)=u^{\operatorname{rk}W}\sum_{j\ge0}u^{-j}c_j(W)$.
The expression is valid also for negative virtual ranks. For $e=1$
and $k=h$, the identity map has a zero-dimensional cokernel and the
sign is $(-1)^b=1$.

Only the coefficient of $u^0$ has the cohomological degree of the
virtual class. Pairing with this class gives the antidiagonal residue
integral of \cite[\S2.2]{bp08}, using the comparison of virtual
classes for the projective target $Y$ \cite[Theorem~0.1]{siebertgw}.
The line degrees $(0,2b)$ sum to $2h-2$, so the genus-variable shift
in the residue convention is zero. Exponentiation of the connected
integrals counts unordered collections of positive-degree connected
components with their automorphism weights. Corollary~7.2 of
\cite{bp08} gives \eqref{eq:positiveseries} for the result.
Transposing the summation partition reverses the content and preserves
the hooks, so either order of the two levels gives this expression.
The same argument on $C$ uses the lines $E$ and $K_C\otimes E$ of
degrees $(0,2c)$. Theorem~\ref{thm:character} now gives the
negative-side series by subtraction in the connected logarithm.
\end{proof}

Write $n_{e,l}$ for the exponent at $(e,l)$ of $P^H(B_1)$
for the character path of Theorem~\ref{thm:character}. Thus $n_{e,l}$ is the jump
of the downstairs exponent in degree $2e$ and genus $l$. For the
balanced factor $A_l$, denote its evaluated logarithm by
\[
 G_l(Q,y)=\log\!\left(\sum_\lambda Q^{|\lambda|}
                  \prod_{a\in\lambda}z_{\ell(a)}^{\,l-1}\right).
\]
This is the balanced specialization \cite[Corollary~7.3]{bp08}.
In particular $[Q]G_l=z^{l-1}$. Finite genus support at each degree
in the geometric factorization gives the coefficientwise identity
\begin{equation}\label{eq:balancedextraction}
 \log Z_b(Q,y)=\sum_{r\ge1}\sum_{l\ge0} n_{r,l}G_l(Q^r,y).
\end{equation}

\begin{theorem}\label{thm:highest}
For the character path, in projection degree $2e$ the highest genus of a changing elementary
exponent is $1+be^2$, and its change is $+1$.
\end{theorem}

\begin{proof}
For a partition of $e$, let $P_R$ and $P_C$ count unordered pairs of
boxes sharing a row and a column, respectively. Counting arms and legs
gives $\sum_a\ell(a)=e+P_R+P_C$ and $c(\lambda)=P_R-P_C$. The
highest power of $y$ contributed by $\lambda$ in
\eqref{eq:positiveseries} is therefore $b(e+2P_R)$. Its maximum
$be^2$ occurs only in the single row and has coefficient $1$.
Products of lower-degree terms in the logarithm have strictly
smaller highest power, since $\sum_i e_i^2<e^2$ for every
decomposition of $e$ into two or more positive summands.

First observe that $n_{r,l}=0$ for $l<1+br$. Riemann--Hurwitz on
the positive-degree components of a connected map of degree $r$ to
$Y$ gives this lower bound on its arithmetic genus. Consequently
$[Q^r]\log Z_b$ vanishes to order at least $2br$ in $v$.
Subtract proper-divisor contributions in \eqref{eq:balancedextraction}
in order of increasing $r$. A preceding factor at degree $d\mid r$
with $l-1\ge bd$ contributes a term of order at least
$2(l-1)r/d\ge 2br$, by the hook formula for $G_l$. The remaining
finite sum $\sum_l n_{r,l}z^{l-1}$ has terms of distinct initial
orders, proving the lower bound by induction.

For $l\ge1$ and $j\ge1$, the highest Laurent power in
$[Q^j]G_l$ is at most $(l-1)j(j+1)/2$. Indeed
$P_R+P_C\le\binom j2$, and products in the logarithm satisfy the
same bound by superadditivity. Suppose inductively that $l-1\le bd^2$
for every nonzero exponent at a smaller degree $d$. At degree $e$
each proper divisor $d$, with $j=e/d>1$, contributes powers no higher
than
\[
 bd^2\,j(j+1)/2 < be^2.
\]
The term $y^{be^2}$ with coefficient $1$ thus survives subtraction
of every proper-divisor term. In the remaining polynomial
$\sum_l n_{e,l}z^{l-1}$ the top Laurent coefficient of $z^{be^2}$
is $(-1)^{be^2}=1$ since $b$ is even. This proves both assertions,
including the starting degree $e=1$.
\end{proof}

\subsection{The quartic correction}
The upstairs degree-one logarithm is $z^b$, so the sole new factor
in projection degree two is $A_{b+1}(q^2)$. The partitions of size
two have hooks $1,2$ and total contents $\pm1$. At upstairs degree
two the quadratic logarithm term $-z^{2b}/2$ is common to
$\log Z_b$ and $G_{b+1}$. Subtracting the latter gives
\begin{equation}\label{eq:quarticgenus}
 \sum_l n_{2,l} z^{l-1}
     = z^{2b}(4-z)^b\big(y^b+y^{-b}-2\big),
\end{equation}
since $z_2=z(4-z)$. The last parenthesis is a polynomial in $z$
with linear term $-b^2 z$, as follows by expansion at $v=0$.
Thus the first changing elementary exponent in projection degree
four occurs in genus $2b+2$, with change $-b^2 4^b$.

For $g=2$ one has $b=2$ and $y^2+y^{-2}-2=-z(4-z)$.
Formula~\eqref{eq:quarticgenus} becomes
\[
 -z^5(4-z)^3=-64z^5+48z^6-12z^7+z^8.
\]
The four jumps occur in genera six through nine. There is already
a connected contribution in genus five. Indeed the leading
coefficient of $[Q^2]\log Z_b$ at $z^{2b}$ is
$2\cdot4^b-\tfrac12=(2^{2h}-1)/2$, which is $63/2$ at $h=3$.
It counts the connected unramified doubles of $Y$ with weight
$1/2$ each and is entirely accounted for by $A_{b+1}(q^2)$.
